% =========================================================================
\section{El amplificador operacional}\label{sec:ao}
% =========================================================================
% -------------------------------------------------------------------------
\subsection{Definición, terminales y circuito equivalente}\label{sec:ao-def}
% -------------------------------------------------------------------------
\begin{definicion}[Amplificador operacional]\label{def:ao}
Un amplificador operacional es un amplificador de tensión con dos entradas, la \emph{inversora}
($-$) y la \emph{no inversora} ($+$), y una salida referida a tierra, que en su zona lineal
amplifica la diferencia de las tensiones de entrada:
\begin{equation}\label{eq:ao}
  \vout=A\,(\vp-\vm),
\end{equation}
donde $A$ es la \emph{ganancia en lazo abierto}. Se alimenta con dos fuentes continuas simétricas,
$+\Vcc$ y $-\Vee$, respecto de la tierra común (en el laboratorio, $\pm10\,\mathrm{V}$).
\end{definicion}

El circuito equivalente de la \figpanel{fig:ao}{b} es el modelo de cuadripolo de la
\cref{sec:cuadripolo} con entrada diferencial: una resistencia de entrada $\Rin$ entre las dos
entradas, y un equivalente de Thévenin $A(\vp-\vm)$ en serie con $\Rout$ a la salida~\cite{TEORIA-EA,SEDRA}.
En el circuito integrado $\mu$A741~\cite{UA741} las patillas son: $2$ (entrada $-$), $3$ (entrada $+$),
$6$ (salida), $7$ ($+\Vcc$), $4$ ($-\Vee$); la $1$ y la $5$ sirven para el ajuste de \emph{offset} y la $8$ no
se conecta. Por simplicidad, los esquemas omiten la alimentación, pero el AO \emph{siempre} está
alimentado: la corriente que entrega a la carga sale de las fuentes, no de la señal.

% Figura: (a) símbolo del AO con alimentación, (b) circuito equivalente, (c) característica en lazo abierto.
\begin{figure}[ht]
\centering
\begin{minipage}[b]{0.26\textwidth}\centering
\begin{circuitikz}[circuito, scale=0.9, transform shape]
  \draw (0,0) node[op amp, ao] (A) {};
  \draw (A.-) to[short, -o] ++(-0.6,0) node[left] {$\vm$};
  \draw (A.+) to[short, -o] ++(-0.6,0) node[left] {$\vp$};
  \draw (A.out) to[short, -o] ++(0.5,0) node[right] {$\vout$};
  \draw (A.up) -- ++(0,0.45) node[above, font=\scriptsize] {$+\Vcc$};
  \draw (A.down) -- ++(0,-0.45) node[below, font=\scriptsize] {$-\Vee$};
  \node[nota] at (-0.95,0.75) {2};
  \node[nota] at (-0.95,-0.75) {3};
  \node[nota] at (0.25,0.75) {7};
  \node[nota] at (0.25,-0.75) {4};
  \node[nota] at (1.15,0.2) {6};
\end{circuitikz}\\[2pt]
{\small (a)}
\end{minipage}\hfill
\begin{minipage}[b]{0.38\textwidth}\centering
\begin{circuitikz}[circuito, scale=0.8, transform shape]
  \draw[azulclaro, fill=opfondo, line width=0.5pt] (-0.4,2.8) -- (-0.4,-0.8) -- (4.9,1.0) -- cycle;
  \draw (-1.1,2.0) node[left] {$\vm$} to[short, o-] (0.6,2.0) to[R, l=$\Rin$] (0.6,0) -- (-1.1,0) to[short, -o] (-1.1,0) node[left] {$\vp$};
  \draw (2.1,-1.1) node[ground] {} -- (2.1,0.1) to[cV, l_=\scriptsize$A(\vp-\vm)$] (2.1,1.9) to[R, l=$\Rout$] (4.4,1.9) to[short, -o] (5.2,1.9) node[right] {$\vout$};
\end{circuitikz}\\[2pt]
{\small (b)}
\end{minipage}\hfill
\begin{minipage}[b]{0.30\textwidth}\centering
\begin{tikzpicture}
  \begin{axis}[temacentro, width=1.15\linewidth, height=4.2cm,
      xmin=-1.6, xmax=1.6, ymin=-1.45, ymax=1.45,
      xtick=\empty, ytick=\empty,
      xlabel={$\vp-\vm$},
      xlabel style={at={(axis description cs:1,0.5)}, anchor=south east, yshift=2pt, font=\scriptsize}]
    \addplot[senal salida, samples=3, domain=-1.6:-0.35] {-1};
    \addplot[senal salida, samples=2, domain=-0.35:0.35] {x/0.35};
    \addplot[senal salida, samples=3, domain=0.35:1.6] {1};
    \addplot[punteado] coordinates {(0.35,0) (0.35,1)};
    \addplot[punteado] coordinates {(-0.35,0) (-0.35,-1)};
    \node[font=\scriptsize, anchor=south] at (axis cs:1.0,1.0) {$+\Vsat$};
    \node[font=\scriptsize, anchor=north] at (axis cs:-1.0,-1.0) {$-\Vsat$};
    \node[font=\scriptsize, anchor=north] at (axis cs:0.42,0.0) {$\frac{\Vsat}{A}$};
    \node[font=\scriptsize, anchor=east, text=rojoteo] at (axis cs:-0.08,0.6) {pendiente $A$};
    \node[font=\scriptsize, anchor=west] at (axis description cs:0.52,0.97) {$\vout$};
  \end{axis}
\end{tikzpicture}\\[2pt]
{\small (c)}
\end{minipage}
\caption{El amplificador operacional. (a) Símbolo, con las patillas del $\mu$A741 y la
alimentación (que en los esquemas siguientes se omite). (b) Circuito equivalente: resistencia
de entrada diferencial $\Rin$ y equivalente de Thévenin de salida~\cite{TEORIA-EA,SEDRA}. (c)
Característica de transferencia en lazo abierto (escala horizontal muy ampliada: la zona lineal
mide $2\Vsat/A\sim\SI{0.1}{\milli\volt}$).}
\label{fig:ao}
\end{figure}


% -------------------------------------------------------------------------
\subsection{Saturación}\label{sec:saturacion}
% -------------------------------------------------------------------------
La salida no puede exceder las tensiones de alimentación. En un AO como el 741 queda, además,
un margen de en torno a $1$--$2\,\mathrm{V}$ con cada fuente~\cite{UA741}:
\begin{equation}\label{eq:saturacion}
  -\Vsat\le\vout\le+\Vsat,\qquad \Vsat\lesssim\Vcc-1\,\mathrm{V}
  \quad(\Vsat\approx9\,\mathrm{V}\ \text{para}\ \Vcc=10\,\mathrm{V}).
\end{equation}
La característica de transferencia en lazo abierto, \cref{eq:ao} con \cref{eq:saturacion}, es
la de la \figpanel{fig:ao}{c}: una recta de pendiente $A$ limitada por $\pm\Vsat$. Como $A\sim10^{5}$,
la zona lineal es solo $\abs{\vp-\vm}<\Vsat/A\sim\SI{90}{\micro\volt}$: sin realimentación, el AO
se comporta como un comparador. Valores concretos del 741 se recogen en la \cref{tab:741}.

% -------------------------------------------------------------------------
\subsection{Modelo ideal y realimentación negativa}\label{sec:ideal}
% -------------------------------------------------------------------------
\begin{hipotesis}[Amplificador operacional ideal]\label{hip:ideal}
\begin{equation}\label{eq:ideal}
  \Rin\to\infty,\qquad \Rout\to0,\qquad A\to\infty\quad(\text{e independiente de la frecuencia}).
\end{equation}
\end{hipotesis}

\begin{proposicion}[Reglas del AO ideal]\label{prop:reglas}
Si el AO ideal trabaja en su zona lineal, gracias a una realimentación negativa (una red que
une la salida con la entrada inversora), se cumplen:
\begin{align}
  &\text{(R1) corrientes de entrada nulas:} && i_{+}=i_{-}=0, \label{eq:R1}\\
  &\text{(R2) cortocircuito virtual:} && \vp=\vm . \label{eq:R2}
\end{align}
Además, por $\Rout=0$, la tensión de salida no depende de la carga conectada.
\end{proposicion}
\begin{proof}
(R1) Entre las entradas solo hay $\Rin$; con $\Rin\to\infty$ no circula corriente por ella.
(R2) En zona lineal, $\vout$ es finita ($\abs{\vout}\le\Vsat$), luego por \cref{eq:ao}
$\vp-\vm=\vout/A\to0$ si $A\to\infty$. La realimentación negativa es la que hace que esta situación
sea estable: si $\vout$ crece, $\vm$ crece con ella, $\vp-\vm$ disminuye y $\vout$ se corrige.
Con la salida en vacío, $\Rout=0$ hace que el divisor de salida de \cref{eq:carga} valga $1$.
\end{proof}

\begin{observacion}[Tierra virtual]\label{obs:tierra-virtual}
Si $\vp=0$ (entrada no inversora a tierra), (R2) da $\vm=0$: la entrada inversora está a potencial
de tierra sin estar conectada a ella. Es una \emph{tierra virtual}: tiene tensión nula pero no absorbe
corriente, por (R1). Las dos reglas son incompatibles con la saturación: si $\abs{\vout}=\Vsat$, el lazo
ya no actúa y deja de cumplirse \cref{eq:R2}.
\end{observacion}

\begin{proposicion}[Método de análisis]\label{prop:metodo}
Todo circuito con AO ideal y realimentación negativa se resuelve con cuatro pasos:
\begin{enumerate}
  \item obtener $\vp$ a partir del circuito conectado a la entrada no inversora (por (R1), sin carga);
  \item imponer $\vm=\vp$ (R2);
  \item aplicar la ley de Kirchhoff de corrientes en el nudo inversor con $i_{-}=0$ (R1), usando
        impedancias si hay elementos reactivos;
  \item despejar $\vout$ y comprobar que $\abs{\vout}<\Vsat$ (si no, el resultado no es válido).
\end{enumerate}
\end{proposicion}

% -------------------------------------------------------------------------
\subsection{Límites del modelo ideal}\label{sec:real}
% -------------------------------------------------------------------------
El modelo ideal es excelente a baja frecuencia y con resistencias de valor moderado
($1$--$100\,\mathrm{k}\Omega$). Sus límites explican lo que se observa al forzar el circuito
(amplitudes grandes o frecuencias de $100\,\mathrm{kHz}$).

\begin{proposicion}[Ganancia en lazo abierto real]\label{prop:Areal}
La ganancia del AO tiene un polo dominante a una frecuencia $f_b$ muy baja~\cite{TEORIA-EA}:
\begin{equation}\label{eq:Areal}
  A(s)=\frac{A_0}{1+s/\omega_b},\qquad
  \abs{A(\jj\omega)}\simeq\frac{A_0\,\omega_b}{\omega}=\frac{\omega_t}{\omega}\quad(\omega\gg\omega_b),\qquad
  f_t=A_0f_b ,
\end{equation}
donde $f_t$, la frecuencia de ganancia unidad, es el \emph{producto ganancia--ancho de banda} (GBW).
\end{proposicion}
\begin{proof}
Es la red STC de la \cref{prop:STC} con $K=A_0$; por encima de $\omega_b$, $\abs{A}=A_0\omega_b/\omega$
se anula ($0\dB$) en $\omega_t=A_0\omega_b$.
\end{proof}

\begin{proposicion}[Ancho de banda en lazo cerrado]\label{prop:BW}
Con ganancia finita $A$, los amplificadores inversor y no inversor de las \cref{sec:pl-inversor,sec:pl-no-inversor} tienen
\begin{equation}\label{eq:lazo-cerrado}
  A_V^{\mathrm{inv}}=\frac{-R_2/R_1}{1+\dfrac{1+R_2/R_1}{A}},\qquad
  A_V^{\mathrm{no\,inv}}=\frac{1+R_2/R_1}{1+\dfrac{1+R_2/R_1}{A}},
\end{equation}
que tienden a los valores ideales si $A\gg1+R_2/R_1$. Con \cref{eq:Areal}, ambos son redes STC paso bajo de frecuencia de corte
\begin{equation}\label{eq:f3dB}
  f_{3\mathrm{dB}}=\frac{f_t}{1+R_2/R_1} .
\end{equation}
\end{proposicion}
\begin{proof}
Con $i_-=0$, el nudo inversor es un divisor entre la entrada y la salida. Para el inversor,
por superposición, $\vm=\vin\,R_2/(R_1+R_2)+\vout\,R_1/(R_1+R_2)$, y $\vout=-A\vm$. Llamando
$\beta=R_1/(R_1+R_2)$, $\vout(1+A\beta)=-A(1-\beta)\vin$, y dividiendo por $A\beta$ se obtiene la
primera expresión, pues $(1-\beta)/\beta=R_2/R_1$ y $1/\beta=1+R_2/R_1$. Para el no inversor,
$\vm=\beta\vout$ y $\vout=A(\vin-\beta\vout)$, de donde $\vout/\vin=(1/\beta)/(1+1/(A\beta))$. Si
$A=\omega_t/s$, el denominador común es $1+s/(\beta\omega_t)$: un polo en $\beta\omega_t$.
\end{proof}

Es decir, ganancia y ancho de banda se intercambian: el producto $\abs{A_V}\cdot f_{3\mathrm{dB}}\approx f_t$ es
constante (\cref{fig:gbw}). Las otras limitaciones relevantes son:
\begin{itemize}
  \item \emph{Velocidad de subida} (\emph{slew rate}): $\abs{\dif\vout/\dif t}\le\mathrm{SR}$. Una sinusoide
        $\vout=V_p\sen\omega t$ tiene pendiente máxima $\omega V_p$; para no deformarse en un
        triángulo debe cumplir
        \begin{equation}\label{eq:slew}
          2\pi f\,V_p\le\mathrm{SR}.
        \end{equation}
  \item \emph{Corriente de salida limitada} (cortocircuito, $\sim\SI{25}{\milli\ampere}$): la salida no
        puede alimentar cargas de pocos ohmios.
  \item \emph{Corrientes de polarización} $I_B$ y \emph{tensión de offset} $V_{IO}$: no nulas; su efecto en
        la salida crece con el valor de las resistencias (o, en un integrador, con el tiempo).
\end{itemize}

% Figura: Bode del AO en lazo abierto (A0=2e5, fb=5 Hz, ft=1 MHz) y de las ganancias en lazo cerrado.
\begin{figure}[ht]
\centering
\begin{tikzpicture}
  \begin{semilogxaxis}[tema, width=0.82\textwidth, height=5.8cm,
      xmin=1, xmax=1e7, ymin=-5, ymax=115, ytick={0,20,40,60,80,100},
      xlabel={$f$ (Hz)}, ylabel={Módulo (dB)}, log basis x=10,
      legend pos=north east]
    \addplot[azulacero, line width=1.1pt, domain=1:1e7, samples=300]
      {20*log10(2e5/sqrt(1+(x/5)^2))};
    \addlegendentry{$\abs{A(\jj\omega)}$, lazo abierto}
    \addplot[rojoteo, line width=1.1pt, domain=1:1e7, samples=300]
      {20*log10(10/sqrt(1+(x/90909)^2))};
    \addlegendentry{$\abs{A_V}$, inversor ($R_2/R_1=10$)}
    \addplot[asintota, domain=1:1e7, samples=2] {20};
    \addplot[punteado] coordinates {(90909,-5) (90909,17)};
    \addplot[punteado] coordinates {(5,-5) (5,103)};
    \addplot[only marks, mark=*, mark size=1.5pt, rojoteo] coordinates {(90909,17)};
    \addplot[only marks, mark=*, mark size=1.5pt, azulacero] coordinates {(5,103) (1e6,0)};
    \node[font=\scriptsize, anchor=south west] at (axis cs:6,103) {$f_b\approx5\,\mathrm{Hz}$};
    \node[font=\scriptsize, anchor=south west] at (axis cs:1.1e6,1) {$f_t\approx1\,\mathrm{MHz}$};
    \node[font=\scriptsize, anchor=west, text=rojoteo] (f3) at (axis cs:150,8) {$f_{3\mathrm{dB}}=f_t/(1+R_2/R_1)\approx91\,\mathrm{kHz}$};
    \draw[->, rojoteo, line width=0.4pt] (f3.east) -- (axis cs:80000,15.5);
    \node[font=\scriptsize, anchor=south, text=grispizarra, rotate=-27] at (axis cs:3000,58) {$-20\dB$/década};
  \end{semilogxaxis}
\end{tikzpicture}
\caption{Diagrama de Bode (módulo) de la ganancia en lazo abierto de un AO de tipo 741, según
\cref{eq:Areal} con $A_0=2\cdot10^{5}$ y $f_t=1\,\mathrm{MHz}$, y de la ganancia en lazo cerrado del
amplificador inversor de la \cref{sec:pl-inversor}, según \cref{eq:lazo-cerrado,eq:f3dB}. La ganancia
en lazo cerrado sigue la recta de $20\dB$ hasta cortar la de lazo abierto: a partir de ahí manda
el AO.}
\label{fig:gbw}
\end{figure}


\begin{table}[ht]
\centering\small
\renewcommand{\arraystretch}{1.2}
\caption{Parámetros típicos del $\mu$A741 a $25\,^{\circ}\mathrm{C}$~\cite{UA741} frente al modelo ideal.}
\label{tab:741}
\begin{tabular}{@{}llcc@{}}
\toprule
Parámetro & Símbolo & $\mu$A741 (típico) & Ideal\\
\midrule
Ganancia en lazo abierto (continua) & $A_0$ & $2\cdot10^{5}$ ($106\dB$) & $\infty$\\
Producto ganancia--ancho de banda & $f_t$ & $\approx\SI{1}{\mega\hertz}$ & $\infty$\\
Resistencia de entrada & $\Rin$ & $\SI{2}{\mega\ohm}$ & $\infty$\\
Resistencia de salida & $\Rout$ & $\SI{75}{\ohm}$ & $0$\\
Velocidad de subida & SR & $\SI{0.5}{\volt\per\micro\second}$ & $\infty$\\
Corriente de polarización & $I_B$ & $\SI{80}{\nano\ampere}$ & $0$\\
Tensión de \emph{offset} de entrada & $V_{IO}$ & $\SI{1}{\milli\volt}$ & $0$\\
Corriente de salida en cortocircuito & $I_{OS}$ & $\pm\SI{25}{\milli\ampere}$ & $\infty$\\
Excursión de salida ($R_L\ge\SI{10}{\kilo\ohm}$) & $\pm\Vsat$ & $\approx\pm(\Vcc-1\,\mathrm{V})$ & $\pm\Vcc$\\
\bottomrule
\end{tabular}
\end{table}
